\chapter{Pembagian}\label{ch:g4:division}

Kelas 3 membagi dengan tabel perkalian dan \omterm{def:g3:division:remainder}{sisa} yang kecil
(\cref{ch:g3:division}); bab ini memasang cara tertulis yang lengkap
--- pembagian bersusun --- yang sanggup menangani bilangan apa pun,
kelompok angka demi kelompok angka.

\section{Cara pembagian bersusun}

\begin{method}[Pembagian bersusun dengan bilangan satu angka]\label{met:g4:division:method}
Untuk membagi $749$ dengan $6$:
\begin{enumerate}
  \item ambil angka paling kiri, $7$ (ratusan): $6$ masuk $1$ kali
        ke dalam $7$; tulis $1$ pada \omterm{def:g3:division:remainder}{hasil bagi}, kurangkan $6$,
        \omterm{def:g3:division:remainder}{sisanya} $1$;
  \item \emph{turunkan} angka berikutnya, $4$: sekarang bagilah $14$
        (puluhan): $6 \times 2 = 12$; tulis $2$, \omterm{def:g3:division:remainder}{sisanya} $2$;
  \item turunkan angka $9$: bagilah $29$: $6 \times 4 = 24$; tulis
        $4$, \omterm{def:g3:division:remainder}{sisanya} $5$;
  \item tidak ada angka yang tersisa: \omterm{def:g3:division:remainder}{hasil baginya} $124$, \omterm{def:g3:division:remainder}{sisanya} $5$.
\end{enumerate}
Periksa: $6 \times 124 + 5 = 744 + 5 = 749$, dan $5 < 6$.
\end{method}

\begin{omfigure}
\begin{tikzpicture}[scale=0.95]
  % kolom angka di x = 0.2, 0.8, 1.4; setiap angka di bawah tempatnya
  \node[font=\normalsize] at (0.2,3.0) {$7$};
  \node[font=\normalsize] at (0.8,3.0) {$4$};
  \node[font=\normalsize] at (1.4,3.0) {$9$};
  % kurung pembagi dan hasil bagi
  \draw[thick] (1.8,2.0) -- (1.8,3.35);
  \draw[thick] (1.8,2.62) -- (3.7,2.62);
  \node[font=\normalsize] at (2.35,3.0) {$6$};
  \node[font=\normalsize, omDef] at (2.35,2.2) {$1$};
  \node[font=\normalsize, omDef] at (2.95,2.2) {$2$};
  \node[font=\normalsize, omDef] at (3.55,2.2) {$4$};
  % langkah 1: 7 - 6 = 1
  \node[font=\normalsize, omThm] at (-0.35,2.4) {$-$};
  \node[font=\normalsize, omThm] at (0.2,2.4) {$6$};
  \draw (-0.55,2.1) -- (0.55,2.1);
  \node[font=\normalsize] at (0.2,1.8) {$1$};
  \node[font=\normalsize] at (0.8,1.8) {$4$};
  \draw[-Stealth, omExo, thin] (0.8,2.75) -- (0.8,2.05);
  % langkah 2: 14 - 12 = 2
  \node[font=\normalsize, omThm] at (-0.35,1.2) {$-$};
  \node[font=\normalsize, omThm] at (0.2,1.2) {$1$};
  \node[font=\normalsize, omThm] at (0.8,1.2) {$2$};
  \draw (-0.55,0.9) -- (1.15,0.9);
  \node[font=\normalsize] at (0.8,0.6) {$2$};
  \node[font=\normalsize] at (1.4,0.6) {$9$};
  \draw[-Stealth, omExo, thin] (1.4,2.75) -- (1.4,0.85);
  % langkah 3: 29 - 24 = 5
  \node[font=\normalsize, omThm] at (0.25,0.0) {$-$};
  \node[font=\normalsize, omThm] at (0.8,0.0) {$2$};
  \node[font=\normalsize, omThm] at (1.4,0.0) {$4$};
  \draw (0.05,-0.3) -- (1.75,-0.3);
  \node[font=\normalsize] at (1.4,-0.6) {$5$};
  % keterangan
  \node[right, font=\small, omExo, align=left] at (4.5,1.4)
    {$7 - 6 = 1$, turunkan angka $4$;\\
     $14 - 12 = 2$, turunkan angka $9$;\\
     $29 - 24 = 5$: berhenti.\\[3pt]
     hasil bagi $124$, sisa $5$\\
     periksa: $6 \times 124 + 5 = 749$};
\end{tikzpicture}

{\small Pembagian bersusun $749$ dengan $6$, ditulis lengkap: setiap
pengurangan tampak, setiap angka tetap di kolomnya, dan anak panah
abu-abu menunjukkan angka yang diturunkan.}
\end{omfigure}

\begin{example}[Nol pada hasil bagi]\label{ex:g4:division:zero}
Bagilah $618$ dengan $3$: ratusan, $6 \div 3 = 2$; puluhan, turunkan
angka $1$: $3$ masuk $0$ kali ke dalam $1$ --- \emph{tulislah} angka
$0$ itu! --- \omterm{def:g3:division:remainder}{sisanya} $1$; satuan, turunkan angka $8$: $18 \div 3 = 6$.
\omterm{def:g3:division:remainder}{Hasil baginya} $206$, \omterm{def:g3:division:remainder}{sisanya} $0$. Melupakan \omterm{def:g1:counting:zero}{nol} di tengah (menulis
$26$) adalah kesalahan yang klasik; pemeriksaan
$3 \times 26 = 78 \neq 618$ langsung menangkapnya.
\end{example}

\begin{example}[Menaksir besarnya hasil bagi]\label{ex:g4:division:size}
Sebelum membagi $749$ dengan $6$, kurunglah dulu jawabannya:
$6 \times 100 = 600$ dan $6 \times 200 = 1\,200$, jadi \omterm{def:g3:division:remainder}{hasil baginya}
antara $100$ dan $200$ --- ia akan punya tiga angka. Taksiran satu
baris ini mencegah sebagian besar kesalahan salah tempat angka.
\end{example}

\section{Memilih apa yang ditanyakan}

\begin{method}[Hasil bagi, sisa, atau hasil bagi ditambah satu]\label{met:g4:division:interpret}
Seperti pada \cref{met:g3:division:interpret}, ceritanya yang menentukan:
\begin{enumerate}
  \item ``berapa \dots yang penuh'' atau ``berapa untuk
        masing-masing'': \omterm{def:g3:division:remainder}{hasil baginya};
  \item ``berapa yang tersisa'': \omterm{def:g3:division:remainder}{sisanya};
  \item ``berapa wadah yang diperlukan untuk semuanya'': \omterm{def:g3:division:remainder}{hasil
        baginya}, ditambah satu kalau \omterm{def:g3:division:remainder}{sisanya} bukan \omterm{def:g1:counting:zero}{nol}.
\end{enumerate}
\end{method}

\begin{example}\label{ex:g4:division:problem}
$530$ buku harus dikemas ke dalam kotak berisi $8$.
\[
530 = 8 \times 66 + 2 .
\]
Berapa kotak yang \emph{penuh}? $66$. Berapa buku yang tidak berada
di kotak yang penuh? $2$. Berapa kotak untuk mengemas \emph{semua}
buku itu? $67$. Tiga pertanyaan, satu pembagian.
\end{example}

\begin{example}[Membagi biaya secara adil]\label{ex:g4:division:money}
Empat orang teman membagi rata biaya sebuah hadiah seharga $92$:
$92 \div 4 = 23$ tepat (\omterm{def:g3:division:remainder}{sisanya} $0$). Masing-masing membayar $23$.
Ketika \omterm{def:g3:division:remainder}{sisanya} \omterm{def:g1:counting:zero}{nol}, pembagiannya ``pas'' dan pembagian biayanya
benar-benar adil.
\end{example}

\section{Latihan}

\begin{exercise}[$\star$]\label{exo:g4:division:1}
Kurunglah setiap \omterm{def:g3:division:remainder}{hasil bagi} seperti pada \cref{ex:g4:division:size}
(di antara dua bilangan ``bulat'' yang mana?), lalu hitunglah
pembagian bersusunnya: $96 \div 4$; \quad $87 \div 5$.
\end{exercise}

\begin{exercise}[$\star$]\label{exo:g4:division:2}
Hitunglah pembagian bersusunnya lalu periksa satu per satu:
$672 \div 4$; \quad $925 \div 7$; \quad $804 \div 6$.
\end{exercise}

\begin{exercise}[$\star$]\label{exo:g4:division:3}
Hati-hati dengan \omterm{def:g1:counting:zero}{nolnya} (lihat \cref{ex:g4:division:zero}):
$816 \div 4$; \quad $420 \div 7$; \quad $2\,512 \div 5$.
\end{exercise}

\begin{exercise}[$\star$]\label{exo:g4:division:4}
Tulislah kesamaan pemeriksaannya ($a = b \times q + r$) untuk:
$85 \div 9$; \quad $1\,000 \div 3$.
\end{exercise}

\begin{exercise}[$\star$]\label{exo:g4:division:5}
$156$ butir telur dimasukkan ke kotak berisi $6$. Berapa kotak yang terisi penuh?
\end{exercise}

\begin{exercise}[$\star$]\label{exo:g4:division:6}
Sehelai pita yang panjangnya $250$ cm dipotong menjadi potongan $8$ cm.
Berapa potongan yang didapat, dan berapa panjang yang tersisa?
\end{exercise}

\begin{exercise}[$\star$]\label{exo:g4:division:7}
$375$ murid pergi menonton pertunjukan dengan bus berisi $52$ tempat
duduk. Berapa bus yang diperlukan? (Ini termasuk keadaan yang mana
pada \cref{met:g4:division:interpret}? Bagilah dengan $52$ memakai
kelipatan: $52 \times 7 = 364$.)
\end{exercise}

\begin{exercise}[$\star$]\label{exo:g4:division:8}
Tiga orang teman membagi rata $234$ kelereng. Berapa kelereng untuk
masing-masing? Periksalah dengan perkalian.
\end{exercise}

\begin{exercise}[$\star$]\label{exo:g4:division:9}
Carilah bilangan yang dibagi: pembagian dengan $7$ memberi \omterm{def:g3:division:remainder}{hasil
bagi} $58$ dan \omterm{def:g3:division:remainder}{sisa} $3$.
\end{exercise}

\begin{exercise}[$\star\star$]\label{exo:g4:division:10}
Seorang pustakawan menata $438$ buku, $9$ buku tiap rak. Rak dijual
dalam lemari berisi $6$ rak. Berapa rak yang diperlukan? Berapa
lemari yang harus dibeli? (Dua pembagian, keduanya dengan
pertanyaan ``ditambah satu''.)
\end{exercise}

\begin{exercise}[$\star\star$]\label{exo:g4:division:11}
Pada sebuah pembagian dengan $8$, \omterm{def:g3:division:remainder}{hasil baginya} sama dengan \omterm{def:g3:division:remainder}{sisanya}.
Bilangan yang dibagi apa saja yang mungkin? Sebutkan semuanya.
(\omterm{def:g3:division:remainder}{Sisanya} harus tetap di bawah $8$.)

\end{exercise}
